%%=============================================================================
%% Conclusie
%%=============================================================================

\chapter{Conclusie}%
\label{ch:conclusie}

% TODO: Trek een duidelijke conclusie, in de vorm van een antwoord op de
% onderzoeksvra(a)g(en). Wat was jouw bijdrage aan het onderzoeksdomein en
% hoe biedt dit meerwaarde aan het vakgebied/doelgroep? 
% Reflecteer kritisch over het resultaat. In Engelse teksten wordt deze sectie
% ``Discussion'' genoemd. Had je deze uitkomst verwacht? Zijn er zaken die nog
% niet duidelijk zijn?
% Heeft het onderzoek geleid tot nieuwe vragen die uitnodigen tot verder 
%onderzoek?


Deze studie toont aan dat neurale netwerken ingezet kunnen worden om de populariteit van een POI te voorspellen op basis van censusdata. De Transformer presteert duidelijk het best. Het behaalt de laagste fouten en de hoogste verklaarde variantie, zowel op normale dagen als op feestdagen. Hoewel de voorspellingsfouten toenemen naarmate de tijd vordert, blijven deze afwijkingen beperkt waardoor de voorspellingen bruikbaar blijven tot minstens een maand in de toekomst.

De resultaten van deze studie kwamen grotendeels overeen met de verwachtingen. Modellen die temporele afhankelijkheden beter kunnen modelleren op zowel korte als lange termijn, zullen ook nauwkeuriger tijdsafhankelijke doelvariabelen kunnen voorspellen dan eenvoudigere modellen die dit niet kunnen. Er wordt toch opgemerkt dat de Transformer moeilijker bezoekersaantallen op feestdagen voorspelt waar afwijken kunnen ontstaan van gebruikelijke bezoekerspatronen dan op reguliere dagen. Daarnaast blijven de voorspellingen ook moeilijk te verklaren omdat een Transformer een black-box model is.

Deze resultaten hebben invloed op verschillende aspecten van de bedrijfswereld. Bedrijven kunnen met dit model betere strategische beslissingen nemen bij een vestiging te openen en gerichter goederen in te kopen. Hierdoor zullen ondernemingen meer omzet kunnen maken en startende bedrijven zullen een grotere kans hebben om te overleven.

Verder kan dit onderzoek leiden tot nieuwe vragen. Hoe nauwkeurig kunnen de bezoekersaantallen voorspeld worden in extreme situaties zoals epidemieën of pandemieën zoals corona? Op welke manieren zou het model hier beter op kunnen voorspellen? Hoe kunnen de voorspellingen eenvoudiger verklaard worden voor bedrijven?
